% !TEX root = ../main.tex

% 结论
\begin{conclusions}

在知识密集型多跳问答任务中，由于模型参数知识的静态性以及复杂推理能力的局限性，
大语言模型面临着知识过时、推理链条不完整以及易产生幻觉等问题。
检索增强生成通过引入外部知识在一定程度上缓解了上述问题，
但在复杂场景下仍存在上下文噪声较大、语义表达粒度单一以及推理能力不足等挑战。
针对上述问题，本文围绕“结构化推理”与“语义建模”两个关键方向展开研究，
分别提出了基于三元组骨架引导的检索增强生成方法以及基于联合语义建模的检索增强生成方法，
以提升多跳问答任务中的知识利用效率与推理能力。本文的主要贡献包括：

（1）针对于复杂多跳问答任务中上下文信噪比较低的问题，本文提出了一种基于三元组骨架引导的检索增强生成方法。
本方法首先从检索文档中抽取知识图谱三元组，并通过迭代扩展构建推理骨架，将其作为结构化推理路径；
在此基础上，通过骨架引导对原始文本进行精确索引，获取与推理路径一致的文本证据，从而构建“结构引导 + 文本补充”的上下文表示方式。
实验结果表明，本方法不仅在各项评价指标上优于对比方法，还能够在保证性能的前提下显著减少输入 Token 长度，
有效降低上下文噪声并提升推理效率与可解释性。

（2）针对于图结构方法依赖关系抽取且构建成本较高的问题，本文提出了一种基于联合语义建模的检索增强生成方法。
本方法在 LinearRAG 框架基础上，对语义表示与上下文构建过程进行了系统优化：
在语义抽取阶段，引入实体与短语的多粒度语义抽取机制，以增强语义覆盖能力；
在表示学习阶段，提出基于句子加权的段落表示方法，更精细地刻画段落内部语义结构；
在上下文构建阶段，提出联合语义上下文范式，将实体、短语与段落信息进行统一组织，从而提升上下文的信息密度与表达能力。
实验结果表明，本方法在无需显式推理路径的情况下，能够实现更稳定且更优的性能表现，在整体指标上优于多种基准方法。

尽管本文针对多跳问答任务中的检索增强生成方法进行了较为系统的研究，并取得了一定的成果，但仍存在一些不足与值得进一步探索的方向：

（1）在方法设计上，分别从结构化推理与语义建模两个角度进行独立优化，尚未充分探索两者的深度融合机制，
未来可研究轻量化推理结构与联合语义建模的统一框架，以进一步提升模型性能。

（2）在语义表示方面，主要依赖静态嵌入模型对实体、短语及句子进行编码，对于语义重要性的建模仍较为粗粒度，
未来可引入可训练的表示学习方法或注意力机制，实现更加精细的语义建模。

（3）在上下文构建策略上，采用固定的检索与组织方式，未来可以探索基于问题复杂度或推理需求的动态上下文构建机制，
以提升系统的灵活性与效率。

\end{conclusions}
